\chapter{Axisymmetric Hankel--Transform Formulation for the Linear pTTM}
\label{app:axisym-hankel-linear-pttm}

This appendix develops a transform-based formulation for the axisymmetric
constant-coefficient pTTM, written in the solver-unit notation of Chapter~\ref{ch:numerical-methods}.
The construction is motivated by the generalized Hankel-transform framework of
Ali and Kalla \cite{Ali_Kalla_1999}, by the classical Hankel-transform
treatment of boundary-value problems in cylindrical geometry
\cite{Davies_2002_Hankel}, and by transform-based analyses of axisymmetric
Dirichlet and heat-conduction problems
\cite{Rogers_2012,Tuteja_Jaloree_Goyal_2014,Kunga_2021}.
The role of the weighted Gelfand--Shilov spaces of type \(W\) developed by
Mahato and Pasawan \cite{Mahato_Pasawan_2021} is auxiliary here: they provide
a convenient transform-invariant test-function setting, but the well-posedness
and convergence arguments below are carried out in the natural weighted
\(L^2\) framework.

\section{Axisymmetric linear pTTM in solver units}
\label{sec:appE-axisym-model}

Let
\begin{equation}
    \xi \coloneqq \frac{\varrho}{L_{\mathrm{ref}}},
    \qquad
    \zeta \coloneqq \frac{z}{L_{\mathrm{ref}}},
    \qquad
    \tau \coloneqq \frac{t_{\mathrm{phys}}}{t_{\mathrm{ref}}},
    \label{eq:appE-scaled-variables}
\end{equation}
where \(L_{\mathrm{ref}}\) and \(t_{\mathrm{ref}}\) are the solver-space
reference scales used in Chapter~\ref{ch:numerical-methods}.
Throughout this appendix, \(\xi\in[0,\infty)\) is the dimensionless radial
    coordinate and \(\zeta\in[0,1]\) is the dimensionless depth coordinate.

    We consider the linear axisymmetric pTTM obtained from the general pTTM of
    Chapter~\ref{ch:theoretical-background} after fixing the material coefficients
    at constant values, suppressing phase-change effects, and subtracting any
    reference or lifting function needed to impose homogeneous boundary data in the
    depth coordinate. The unknowns are denoted by
    \[
        u(\tau,\xi,\zeta), \qquad v(\tau,\xi,\zeta),
    \]
    representing the scaled electron and lattice temperature fields, respectively.
    They satisfy
    \begin{subequations}\label{eq:appE-linear-axisym-pttm}
        \begin{align}
            \partial_\tau u
             & =
            D_e
            \left(
            \partial_{\xi\xi}u + \frac{1}{\xi}\partial_\xi u + \partial_{\zeta\zeta}u
            \right)
            - c_e (u-v)
            + s(\tau,\xi,\zeta),
            \label{eq:appE-linear-axisym-pttm-e}
            \\
            \partial_\tau v
             & =
            D_l
            \left(
            \partial_{\xi\xi}v + \frac{1}{\xi}\partial_\xi v + \partial_{\zeta\zeta}v
            \right)
            + c_l (u-v),
            \label{eq:appE-linear-axisym-pttm-l}
        \end{align}
    \end{subequations}
    for \((\tau,\xi,\zeta)\in(0,\tau_f]\times(0,\infty)\times(0,1)\), where
\(D_e,D_l>0\) are constant diffusivities and
\(c_e,c_l>0\) are constant coupling parameters.

The radial operator is written in self-adjoint axisymmetric form as
\begin{equation}
    \mathcal{L}_\xi w
    \coloneqq
    \partial_{\xi\xi}w + \frac{1}{\xi}\partial_\xi w
    =
    \frac{1}{\xi}\partial_\xi\!\bigl(\xi\,\partial_\xi w\bigr).
    \label{eq:appE-radial-operator}
\end{equation}

At the symmetry axis \(\xi=0\) we impose the regularity conditions
\begin{equation}
    \partial_\xi u(\tau,0,\zeta)=0,
    \qquad
    \partial_\xi v(\tau,0,\zeta)=0,
    \label{eq:appE-axis-regularity}
\end{equation}
and as \(\xi\to\infty\) we assume sufficient decay so that the Hankel
integrals and integration-by-parts identities below are valid.

For definiteness, this appendix is written for homogeneous Dirichlet conditions
in the depth variable,
\begin{equation}
    u(\tau,\xi,0)=u(\tau,\xi,1)=0,
    \qquad
    v(\tau,\xi,0)=v(\tau,\xi,1)=0.
    \label{eq:appE-depth-dirichlet}
\end{equation}
The homogeneous Neumann case is stated separately in
Remark~\ref{rem:appE-neumann-case}; in that case the depth basis becomes the
cosine basis, which is the setting closest to a discrete cosine transform (DCT) implementation
\cite{Rogers_2012}.

The initial conditions are
\begin{equation}
    u(0,\xi,\zeta)=u_0(\xi,\zeta),
    \qquad
    v(0,\xi,\zeta)=v_0(\xi,\zeta).
    \label{eq:appE-initial-data}
\end{equation}

\section{Generalized Hankel transform and specialisation to the axisymmetric kernel}
\label{sec:appE-generalized-hankel}

Ali and Kalla \cite{Ali_Kalla_1999} introduce the generalized Hankel transform
\begin{equation}
    \mathscr{H}_{\nu}^{a,c}[f](s)
    \coloneqq
    \int_{0}^{\infty}
    \xi^{a} J_{\nu}(s\xi^{c})\, f(\xi)\, d\xi,
    \qquad c>0,
    \label{eq:appE-generalized-hankel}
\end{equation}
with kernel
\begin{equation}
    K_{\nu}^{a,c}(s,\xi)\coloneqq \xi^{a}J_{\nu}(s\xi^{c}),
    \label{eq:appE-generalized-hankel-kernel}
\end{equation}
and corresponding inversion
\begin{equation}
    f(\xi)
    =
    c\,\xi^{2c-a-1}
    \int_{0}^{\infty}
    s\, \mathscr{H}_{\nu}^{a,c}[f](s)\,
    J_{\nu}(s\xi^{c})\, ds.
    \label{eq:appE-generalized-hankel-inverse}
\end{equation}
The classical Hankel transform of order \(\nu\) is recovered by setting
\(a=c=1\) \cite{Ali_Kalla_1999,Davies_2002_Hankel}.

For the axisymmetric linear pTTM, the relevant specialisation is the
order-zero classical transform
\begin{equation}
    \mathscr{H}_0[f](\eta)
    \coloneqq
    \int_{0}^{\infty}
    \xi f(\xi) J_0(\eta\xi)\, d\xi,
    \qquad \eta\ge 0,
    \label{eq:appE-hankel-order-zero}
\end{equation}
which corresponds to the parameter choice
\((\nu,a,c)=(0,1,1)\) in \eqref{eq:appE-generalized-hankel}.
Its inverse is
\begin{equation}
    f(\xi)
    =
    \int_{0}^{\infty}
    \eta\, \mathscr{H}_0[f](\eta)\, J_0(\eta\xi)\, d\eta,
    \label{eq:appE-hankel-order-zero-inverse}
\end{equation}
and Parseval's identity takes the form
\begin{equation}
    \int_{0}^{\infty}\xi f(\xi)g(\xi)\, d\xi
    =
    \int_{0}^{\infty}\eta\, \mathscr{H}_0[f](\eta)\, \mathscr{H}_0[g](\eta)\, d\eta.
    \label{eq:appE-hankel-parseval}
\end{equation}

The central operational identity is the diagonalisation of the radial operator.

\begin{proposition}
    \label{prop:appE-hankel-diagonalisation}
    Let $f \in C^2((0,\infty))$
    and let $\mathcal{L}_\xi$ denote the radial part of the Laplacian operator defined in \eqref{eq:appE-radial-operator}.
    Assume that $f, f', f''$ are such that the Hankel transform of order zero of \eqref{eq:appE-hankel-order-zero}
    exists and converges absolutely for $\eta > 0$. If $f$ satisfies the following asymptotic kernel-weighted boundary conditions:
    \begin{enumerate}
        \item $\lim_{\xi \to \{0, \infty\}} \xi f_\xi(\xi) J_0(\eta\xi) = 0$
        \item $\lim_{\xi \to \{0, \infty\}} \eta \xi f(\xi) J_1(\eta\xi) = 0$
    \end{enumerate}
    then the operator $\mathcal{L}_\xi$ is diagonalised by $\mathscr{H}_0$ such that:
    \begin{equation}
        \mathscr{H}_0[\mathcal{L}_\xi f](\eta) = -\eta^2 \mathscr{H}_0[f](\eta).
        \label{eq:appE-hankel-radial-diag}
    \end{equation}
\end{proposition}

\begin{proof}
    Using the divergence form \eqref{eq:appE-radial-operator},
    \[
        \mathscr{H}_0[\mathcal{L}_\xi f](\eta)
        =
        \int_0^\infty
        \partial_\xi(\xi f_\xi)\, J_0(\eta\xi)\, d\xi.
    \]
    Integrating by parts once yields
    \[
        \mathscr{H}_0[\mathcal{L}_\xi f](\eta)
        =
        -\int_0^\infty
        \xi f_\xi(\xi)\, \partial_\xi J_0(\eta\xi)\, d\xi.
    \]
    Integrating by parts again gives
    \[
        \mathscr{H}_0[\mathcal{L}_\xi f](\eta)
        =
        \int_0^\infty
        f(\xi)\,
        \partial_\xi\!\bigl(\xi\,\partial_\xi J_0(\eta\xi)\bigr)\, d\xi.
    \]
    The Bessel equation for \(J_0\) implies
    \[
        \partial_{\xi\xi}J_0(\eta\xi) + \frac{1}{\xi}\partial_\xi J_0(\eta\xi)
        + \eta^2 J_0(\eta\xi)=0,
    \]
    hence
    \[
        \partial_\xi\!\bigl(\xi\,\partial_\xi J_0(\eta\xi)\bigr)
        =
        -\eta^2 \xi J_0(\eta\xi).
    \]
    Substituting this identity proves \eqref{eq:appE-hankel-radial-diag}.
\end{proof}

\begin{remark}
    The generalized transform of Ali and Kalla \cite{Ali_Kalla_1999} is useful for
    placing the construction in a broader kernel framework, but for the
    axisymmetric linear pTTM the order-zero classical Hankel transform already
    diagonalises the radial Bessel operator. This is the basic reason that
    the axisymmetric transform machinery is especially natural in this setting
    \cite{Davies_2002_Hankel,Rogers_2012}.
\end{remark}

\section{Depth expansion and modal reduction}
\label{sec:appE-depth-expansion}

Under the homogeneous Dirichlet conditions
\eqref{eq:appE-depth-dirichlet}, the natural depth eigenfunctions are
\begin{equation}
    \phi_n(\zeta)\coloneqq \sqrt{2}\sin(n\pi\zeta),
    \qquad
    \mu_n\coloneqq n\pi,
    \qquad
    n\in\mathbb{N},
    \label{eq:appE-depth-basis}
\end{equation}
which satisfy
\[
    -\phi_n''=\mu_n^2\phi_n,
    \qquad
    \phi_n(0)=\phi_n(1)=0,
    \qquad
    \int_0^1 \phi_n(\zeta)\phi_m(\zeta)\, d\zeta=\delta_{nm}.
\]

We expand
\begin{equation}
    u(\tau,\xi,\zeta)
    =
    \sum_{n=1}^{\infty} u_n(\tau,\xi)\phi_n(\zeta),
    \qquad
    v(\tau,\xi,\zeta)
    =
    \sum_{n=1}^{\infty} v_n(\tau,\xi)\phi_n(\zeta),
    \label{eq:appE-depth-expansion}
\end{equation}
with modal source coefficients
\begin{equation}
    s_n(\tau,\xi)
    \coloneqq
    \int_0^1 s(\tau,\xi,\zeta)\phi_n(\zeta)\, d\zeta.
    \label{eq:appE-source-depth-coeff}
\end{equation}

Applying the order-zero Hankel transform in \(\xi\), define
\begin{equation}
    \widehat{u}_n(\tau,\eta)
    \coloneqq
    \int_0^\infty \xi u_n(\tau,\xi)J_0(\eta\xi)\, d\xi,
    \qquad
    \widehat{v}_n(\tau,\eta)
    \coloneqq
    \int_0^\infty \xi v_n(\tau,\xi)J_0(\eta\xi)\, d\xi,
    \label{eq:appE-modal-hankel-coeffs}
\end{equation}
and
\begin{equation}
    \widehat{s}_n(\tau,\eta)
    \coloneqq
    \int_0^\infty \xi s_n(\tau,\xi)J_0(\eta\xi)\, d\xi.
    \label{eq:appE-modal-source-hankel}
\end{equation}

By Proposition~\ref{prop:appE-hankel-diagonalisation},
each transformed mode satisfies a \(2\times 2\) linear ODE:
\begin{equation}
    \partial_\tau
    \widehat{\mathbf{w}}_n(\tau,\eta)
    =
    \mathbf{A}_n(\eta)\widehat{\mathbf{w}}_n(\tau,\eta)
    +
    \widehat{s}_n(\tau,\eta)\mathbf{e}_1,
    \qquad
    \mathbf{e}_1\coloneqq (1,0)^\top,
    \label{eq:appE-modal-ode}
\end{equation}
where
\begin{equation}
    \widehat{\mathbf{w}}_n
    \coloneqq
    \begin{bmatrix}
        \widehat{u}_n \\[2pt]
        \widehat{v}_n
    \end{bmatrix},
    \qquad
    \mathbf{A}_n(\eta)
    \coloneqq
    \begin{bmatrix}
        -D_e\omega_n(\eta)-c_e & c_e                    \\
        c_l                    & -D_l\omega_n(\eta)-c_l
    \end{bmatrix},
    \label{eq:appE-modal-matrix}
\end{equation}
with
\begin{equation}
    \omega_n(\eta)\coloneqq \eta^2+\mu_n^2.
    \label{eq:appE-omega-def}
\end{equation}

Hence the original axisymmetric PDE system is reduced to a doubly indexed
family of linear ODE systems in \((\tau,\eta,n)\), as in classical
Hankel/Fourier transform methods for axisymmetric heat-type equations
\cite{Davies_2002_Hankel,Rogers_2012,Kunga_2021}.

\section{Modal solution formula and spectral stability}
\label{sec:appE-modal-solution}

The solution of \eqref{eq:appE-modal-ode} is given by Duhamel's formula:\footnote{For a comprehensive account of the Duhamel integral and its various applications, see R{\'o}{\.z}a{\'n}ski et al. \cite{Różański_2021}.}
\begin{equation}
    \widehat{\mathbf{w}}_n(\tau,\eta)
    =
    e^{\tau \mathbf{A}_n(\eta)}
    \widehat{\mathbf{w}}_n(0,\eta)
    +
    \int_0^\tau
    e^{(\tau-\sigma)\mathbf{A}_n(\eta)}
    \widehat{s}_n(\sigma,\eta)\mathbf{e}_1\, d\sigma.
    \label{eq:appE-duhamel}
\end{equation}

The modal matrix has trace
\begin{equation}
    \operatorname{tr}\mathbf{A}_n(\eta)
    =
    -(D_e+D_l)\omega_n(\eta)-(c_e+c_l)<0,
    \label{eq:appE-trace}
\end{equation}
and determinant
\begin{equation}
    \det \mathbf{A}_n(\eta)
    =
    D_eD_l\,\omega_n(\eta)^2
    +
    (D_ec_l + D_lc_e)\omega_n(\eta)
    \ge 0.
    \label{eq:appE-determinant}
\end{equation}
Therefore both eigenvalues of \(\mathbf{A}_n(\eta)\) have non-positive real
part.

\begin{proposition}
    \label{prop:appE-modal-stability}
    For every \(\eta\ge 0\) and \(n\ge 1\), the matrix
    \(\mathbf{A}_n(\eta)\) generates a contractive semigroup on \(\mathbb{R}^2\).\footnote{
    In this context, the semigroup $\{e^{\tau \mathbf{A}_n(\eta)}\}_{\tau \geq 0}$ is contractive if $\|e^{\tau \mathbf{A}_n(\eta)}\| \leq 1$ for all $\tau \geq 0$ with respect to a suitable norm.
    } In particular, there exists a constant \(C\ge 1\), independent of
    \((\eta,n)\), such that
    \begin{equation}
        \left\|
        e^{\tau \mathbf{A}_n(\eta)}
        \right\|
        \le
        C\, e^{-c_\star \tau},
        \qquad
        c_\star \coloneqq \frac{1}{2}\min\{c_e,c_l\},
        \qquad
        \tau\ge 0.
        \label{eq:appE-modal-semigroup-bound}
    \end{equation}
\end{proposition}

\begin{proof}
    The claim follows from \eqref{eq:appE-trace} and
    \eqref{eq:appE-determinant}, together with the fact that the symmetric part of
    \(\mathbf{A}_n(\eta)\) is strictly negative definite up to the coupling
    weights. Since \(\omega_n(\eta)\ge \pi^2\), the diffusive contribution adds
    further damping and cannot destabilize the system.
\end{proof}

The inverse reconstruction is
\begin{equation}
    \mathbf{w}(\tau,\xi,\zeta)
    =
    \sum_{n=1}^{\infty}
    \phi_n(\zeta)
    \int_0^\infty
    \eta\, J_0(\eta\xi)\,
    \widehat{\mathbf{w}}_n(\tau,\eta)\, d\eta,
    \qquad
    \mathbf{w}\coloneqq (u,v)^\top.
    \label{eq:appE-inverse-reconstruction}
\end{equation}

\section{Energy estimate and well-posedness}
\label{sec:appE-wellposedness}

Define the weighted Lebesgue space
\begin{equation}
    \mathcal{X} \coloneqq L^2\bigl((0,\infty)\times(0,1);\, \xi\, d\xi\, d\zeta\bigr),
    \label{eq:appE-weighted-space}
\end{equation}
equipped with the inner product
\begin{equation}
    \langle f,g\rangle_\mathcal{X} \coloneqq \int_0^1\!\!\int_0^\infty f(\xi,\zeta)g(\xi,\zeta)\,\xi\, d\xi\, d\zeta.
    \label{eq:appE-weighted-inner-product}
\end{equation}
To accommodate the diffusion operators, we introduce the weighted Sobolev space
\begin{equation}
    \mathcal{V} \coloneqq \bigl\{ \phi \in \mathcal{X} : \partial_\xi \phi \in \mathcal{X}, \, \partial_\zeta \phi \in \mathcal{X}, \, \phi|_{\zeta \in \{0,1\}} = 0 \bigr\},
\end{equation}
where derivatives are understood in the distributional sense.
The natural energy space for the coupled system is the product space $\mathbb{X} \coloneqq \mathcal{X} \times \mathcal{X}$.

\begin{theorem}
    \label{thm:appE-energy-estimate}
    Let $(u_0,v_0)\in \mathbb X$ and let $s\in L^2(0,\tau_f;\mathcal X)$.
    Then the system \eqref{eq:appE-linear-axisym-pttm}--\eqref{eq:appE-initial-data}
    admits a unique mild solution
    \[
        (u,v)\in C([0,\tau_f];\mathbb X).
    \]
    If, in addition, $(u_0,v_0)\in D(\mathcal A)$, where
    \[
        D(\mathcal A)
        =
        \Bigl\{(u,v)\in \mathcal V\times\mathcal V:
        \mathcal L_\xi u+\partial_{\zeta\zeta}u\in\mathcal X,\
        \mathcal L_\xi v+\partial_{\zeta\zeta}v\in\mathcal X
        \Bigr\},
    \]
    then $(u,v)$ is a strong solution and, for a.e. $\tau\in(0,\tau_f)$, satisfies
    \begin{align}
        \frac{d}{d\tau}\mathcal E(\tau)
        + D_e\|\nabla u\|_\mathcal X^2
        + \frac{c_e}{c_l}D_l\|\nabla v\|_\mathcal X^2
        + c_e\|u-v\|_\mathcal X^2
        =
        \langle s(\tau),u(\tau)\rangle_\mathcal X,
        \label{eq:appE-energy-identity}
    \end{align}
    with
    \begin{equation}
        \mathcal{E}(\tau) \coloneqq \frac{1}{2}\|u(\tau)\|_\mathcal{X}^2 + \frac{c_e}{2c_l}\|v(\tau)\|_\mathcal{X}^2.
        \label{eq:appE-energy-functional}
    \end{equation}
    Consequently,
    \begin{equation}
        \sup_{0\le \tau\le \tau_f}\mathcal E(\tau)
        \le
        e^{\tau_f}
        \left(
        \mathcal E(0)+\|s\|_{L^2(0,\tau_f;\mathcal X)}^2
        \right).
        \label{eq:appE-energy-bound}
    \end{equation}
\end{theorem}

\begin{proof}
    We divide the argument into two steps.

    \smallskip
    \noindent
    \textit{Step 1: existence and uniqueness of the mild solution.}
    Let
    \[
        \phi_n(\zeta) \coloneqq \sqrt{2}\sin(n\pi\zeta),
        \qquad n\in\mathbb{N},
    \]
    be the orthonormal Dirichlet basis in \(L^2(0,1)\), and let
    \(\mathscr H_0\) denote the order-zero Hankel transform in the radial
    variable \(\xi\). Expanding \(u\), \(v\), and \(s\) in the depth basis and
    then applying \(\mathscr H_0\) in \(\xi\), one obtains for each
    \((n,\eta)\in \mathbb{N}\times [0,\infty)\) the transformed modal system
    \[
        \partial_\tau \widehat{\mathbf w}_n(\tau,\eta)
        =
        \mathbf A_n(\eta)\widehat{\mathbf w}_n(\tau,\eta)
        +
        \widehat{s}_n(\tau,\eta)\mathbf e_1,
        \qquad
        \widehat{\mathbf w}_n(0,\eta)=\widehat{\mathbf w}_{n,0}(\eta),
    \]
    where
    \[
        \widehat{\mathbf w}_n
        =
        \begin{bmatrix}
            \widehat u_n \\[2pt]
            \widehat v_n
        \end{bmatrix},
        \qquad
        \mathbf e_1=
        \begin{bmatrix}
            1 \\[2pt]
            0
        \end{bmatrix},
    \]
    and \(\mathbf A_n(\eta)\) is the \(2\times 2\) modal matrix introduced in
    \eqref{eq:appE-modal-matrix}. By
    Proposition~\ref{prop:appE-modal-stability}, the family
    \(e^{\tau \mathbf A_n(\eta)}\) is uniformly bounded for
    \(\tau\in[0,\tau_f]\). Hence, for each \((n,\eta)\), Duhamel's formula gives
    the unique modal solution
    \[
        \widehat{\mathbf w}_n(\tau,\eta)
        =
        e^{\tau \mathbf A_n(\eta)}\widehat{\mathbf w}_{n,0}(\eta)
        +
        \int_0^\tau
        e^{(\tau-\sigma)\mathbf A_n(\eta)}
        \widehat{s}_n(\sigma,\eta)\mathbf e_1\, d\sigma.
    \]

    Using Parseval's identity for \(\mathscr H_0\),
    equation~\eqref{eq:appE-hankel-parseval}, and the orthonormality of
    \(\{\phi_n\}_{n\ge 1}\), one reconstructs from these modal solutions a unique
    pair $(u,v)\in C([0,\tau_f];\mathbb X)$, which proves the existence and uniqueness of the mild solution.

    \smallskip
    \noindent
    \textit{Step 2: energy identity and a priori estimate.}
    Assume now that \((u_0,v_0)\in D(\mathcal A)\), so that \((u,v)\) is a strong
    solution. To justify the weighted integration by parts, one may
    either work directly at the level of strong solutions or
    establish the identity first for Galerkin approximants in the modal basis and
    then pass to the limit. We proceed formally, noting that the Galerkin passage
    is standard and yields the same identity.

    Test \eqref{eq:appE-linear-axisym-pttm-e} with \(u\) and
    \eqref{eq:appE-linear-axisym-pttm-l} with \(\frac{c_e}{c_l}v\) in
    \(\mathcal X\). For a sufficiently regular function \(\phi\) satisfying the
    Dirichlet conditions in \(\zeta\), one has
    \begin{align*}
        \bigl\langle \mathcal L_\xi \phi,\phi\bigr\rangle_{\mathcal X}
         & =
        \int_0^1\!\!\int_0^\infty
        \frac{1}{\xi}\partial_\xi(\xi\partial_\xi\phi)\,\phi\,
        \xi\, d\xi\, d\zeta
        =
        \int_0^1\!\!\int_0^\infty
        \partial_\xi(\xi\partial_\xi\phi)\,\phi\, d\xi\, d\zeta
        \\
         & =
        -\int_0^1\!\!\int_0^\infty
        \xi\, |\partial_\xi\phi|^2\, d\xi\, d\zeta
        =
        -\|\partial_\xi\phi\|_{\mathcal X}^2,
    \end{align*}
    where the boundary terms vanish at \(\xi=0\) by axis regularity, at
    \(\xi\to\infty\) by the decay inherited from the Galerkin approximation (and
    then by passage to the limit), and at \(\zeta\in\{0,1\}\) by the homogeneous
    Dirichlet condition. Similarly,
    \[
        \bigl\langle \partial_{\zeta\zeta}\phi,\phi\bigr\rangle_{\mathcal X}
        =
        -\|\partial_\zeta\phi\|_{\mathcal X}^2.
    \]
    Therefore,
    \[
        \bigl\langle \mathcal L_\xi\phi+\partial_{\zeta\zeta}\phi,\phi
        \bigr\rangle_{\mathcal X}
        =
        -\|\partial_\xi\phi\|_{\mathcal X}^2
        -\|\partial_\zeta\phi\|_{\mathcal X}^2
        =
        -\|\nabla \phi\|_{\mathcal X}^2.
    \]

    Applying this with \(\phi=u\) and \(\phi=v\), and summing the two tested
    equations, yields
    \begin{align*}
        \frac{1}{2}\frac{d}{d\tau}\|u(\tau)\|_{\mathcal X}^2
        +
        \frac{c_e}{2c_l}\frac{d}{d\tau}\|v(\tau)\|_{\mathcal X}^2
         & +
        D_e\|\nabla u(\tau)\|_{\mathcal X}^2
        +
        \frac{c_e}{c_l}D_l\|\nabla v(\tau)\|_{\mathcal X}^2
        \\
         & \qquad
        -c_e\langle u-v,u\rangle_{\mathcal X}
        +
        c_e\langle u-v,v\rangle_{\mathcal X}
        =
        \langle s(\tau),u(\tau)\rangle_{\mathcal X}.
    \end{align*}
    The coupling terms combine as
    \[
        -c_e\langle u-v,u\rangle_{\mathcal X}
        +
        c_e\langle u-v,v\rangle_{\mathcal X}
        =
        -c_e\|u-v\|_{\mathcal X}^2.
    \]
    Hence
    \eqref{eq:appE-energy-identity} follows, with
    \(\mathcal E(\tau)\) given by \eqref{eq:appE-energy-functional}.

    To obtain the stability estimate, discard the non-negative dissipation terms in
    \eqref{eq:appE-energy-identity}. Then
    \[
        \frac{d}{d\tau}\mathcal E(\tau)
        \le
        \langle s(\tau),u(\tau)\rangle_{\mathcal X}.
    \]
    By Cauchy--Schwarz and Young's inequality,
    \[
        \langle s(\tau),u(\tau)\rangle_{\mathcal X}
        \le
        \|s(\tau)\|_{\mathcal X}\|u(\tau)\|_{\mathcal X}
        \le
        \frac12\|s(\tau)\|_{\mathcal X}^2
        +
        \frac12\|u(\tau)\|_{\mathcal X}^2.
    \]
    Since
    \[
        \mathcal E(\tau)
        =
        \frac12\|u(\tau)\|_{\mathcal X}^2
        +
        \frac{c_e}{2c_l}\|v(\tau)\|_{\mathcal X}^2
        \ge
        \frac12\|u(\tau)\|_{\mathcal X}^2,
    \]
    it follows that
    \[
        \frac{d}{d\tau}\mathcal E(\tau)
        \le
        \mathcal E(\tau)
        +
        \frac12\|s(\tau)\|_{\mathcal X}^2.
    \]
    Integrating over \((0,\tau)\) gives
    \[
        \mathcal E(\tau)
        \le
        \mathcal E(0)
        +
        \int_0^\tau \mathcal E(t)\, dt
        +
        \frac12\int_0^\tau \|s(t)\|_{\mathcal X}^2\, dt.
    \]
    An application of Gr\"onwall's inequality yields
    \[
        \mathcal E(\tau)
        \le
        e^\tau
        \left(
        \mathcal E(0)
        +
        \frac12\int_0^\tau \|s(t)\|_{\mathcal X}^2\, dt
        \right)
        \le
        e^{\tau_f}
        \left(
        \mathcal E(0)
        +
        \|s\|_{L^2(0,\tau_f;\mathcal X)}^2
        \right).
    \]
    Taking the supremum over \(0\le \tau\le \tau_f\) proves
    \eqref{eq:appE-energy-bound}.
\end{proof}

\begin{remark}
    If one wishes to begin on a transform-invariant dense test-function space
    before passing to \(\mathcal{X}\), the weighted Gelfand--Shilov framework of type \(W\)
    developed for Hankel-related operators in \cite{Mahato_Pasawan_2021} provides
    a natural alternative. In the present appendix, however, the weighted
    \(L^2\)-theory is already sufficient for both the modal representation and the
    convergence analysis.
\end{remark}

\section{Convergence of truncated Hankel--eigenfunction expansions}
\label{sec:appE-convergence}

For \(K>0\) and \(N\in\mathbb{N}\), define the truncated reconstruction
\begin{equation}
    \mathbf{w}_{K,N}(\tau,\xi,\zeta)
    \coloneqq
    \sum_{n=1}^{N}
    \phi_n(\zeta)
    \int_0^{K}
    \eta\, J_0(\eta\xi)\,
    \widehat{\mathbf{w}}_n(\tau,\eta)\, d\eta.
    \label{eq:appE-truncated-reconstruction}
\end{equation}

\begin{theorem}
    \label{thm:appE-spectral-convergence}
    Let \((u,v)\) be the solution of
    \eqref{eq:appE-linear-axisym-pttm}--\eqref{eq:appE-initial-data} with
    \((u_0,v_0)\in \mathcal{X}\times \mathcal{X}\) and \(s\in L^2(0,\tau_f;\mathcal{X})\).
    Then
    \begin{equation}
        \sup_{0\le \tau\le \tau_f}
        \left\|
        \mathbf{w}(\tau)-\mathbf{w}_{K,N}(\tau)
        \right\|_{\mathcal{X}\times \mathcal{X}}
        \longrightarrow 0
        \qquad
        \text{as } K\to\infty,\; N\to\infty.
        \label{eq:appE-spectral-convergence}
    \end{equation}
    More precisely,
    \begin{equation}
        \sup_{0\le \tau\le \tau_f}
        \left\|
        \mathbf{w}(\tau)-\mathbf{w}_{K,N}(\tau)
        \right\|_{\mathcal{X}\times \mathcal{X}}^2
        \le
        C(\tau_f)\Bigl(\mathcal{T}_0(K,N)+\mathcal{T}_s(K,N)\Bigr),
        \label{eq:appE-spectral-tail-bound}
    \end{equation}
    where
    \begin{align}
        \mathcal{T}_0(K,N)
         & \coloneqq
        \sum_{n>N}\int_0^\infty
        \eta\,\|\widehat{\mathbf{w}}_n(0,\eta)\|_{\mathbb{R}^2}^2\, d\eta
        +
        \sum_{n=1}^{N}\int_K^\infty
        \eta\,\|\widehat{\mathbf{w}}_n(0,\eta)\|_{\mathbb{R}^2}^2\, d\eta,
        \label{eq:appE-tail-initial}
        \\
        \mathcal{T}_s(K,N)
         & \coloneqq
        \int_0^{\tau_f}
        \left[
            \sum_{n>N}\int_0^\infty \eta\, |\widehat{s}_n(\tau,\eta)|^2\, d\eta
            +
            \sum_{n=1}^{N}\int_K^\infty \eta\, |\widehat{s}_n(\tau,\eta)|^2\, d\eta
            \right] d\tau.
        \label{eq:appE-tail-source}
    \end{align}
\end{theorem}

\begin{proof}
    Let
    \[
        \mathbf w(\tau,\xi,\zeta)\coloneqq
        \begin{bmatrix}
            u(\tau,\xi,\zeta) \\[2pt]
            v(\tau,\xi,\zeta)
        \end{bmatrix},
        \qquad
        \widehat{\mathbf w}_n(\tau,\eta)\coloneqq
        \begin{bmatrix}
            \widehat u_n(\tau,\eta) \\[2pt]
            \widehat v_n(\tau,\eta)
        \end{bmatrix},
    \]
    where \(\{\phi_n\}_{n\ge 1}\) is the depth eigenbasis and
    \(\widehat{\mathbf w}_n\) denotes the order-zero Hankel transform of the
    \(n\)-th depth mode. We also write
    \[
        \mathbf e_1 \coloneqq
        \begin{bmatrix}
            1 \\[2pt]
            0
        \end{bmatrix}.
    \]

    For each \(n\in\mathbb N\) and \(\eta\ge 0\), the modal coefficients satisfy
    the transformed system
    \begin{equation}
        \partial_\tau \widehat{\mathbf w}_n(\tau,\eta)
        =
        \mathbf A_n(\eta)\widehat{\mathbf w}_n(\tau,\eta)
        +
        \widehat s_n(\tau,\eta)\mathbf e_1,
        \qquad
        \widehat{\mathbf w}_n(0,\eta)=\widehat{\mathbf w}_{n,0}(\eta),
        \label{eq:appE-proof-modal-ode}
    \end{equation}
    hence, by Duhamel's formula \eqref{eq:appE-duhamel},
    \begin{equation}
        \widehat{\mathbf w}_n(\tau,\eta)
        =
        e^{\tau \mathbf A_n(\eta)}\widehat{\mathbf w}_{n,0}(\eta)
        +
        \int_0^\tau
        e^{(\tau-\sigma)\mathbf A_n(\eta)}
        \widehat s_n(\sigma,\eta)\mathbf e_1\, d\sigma.
        \label{eq:appE-proof-duhamel}
    \end{equation}

    \smallskip
    \noindent
    \textit{Step 1: exact representation of the truncation error.}
    Since \(\{\phi_n\}_{n\ge 1}\) is orthonormal in \(L^2(0,1)\) and the
    order-zero Hankel transform is an isometry on
    \(L^2((0,\infty);\xi\,d\xi)\), Parseval's identity yields, for every
    \(\tau\in[0,\tau_f]\),
    \begin{align}
        \bigl\|\mathbf w(\tau)-\mathbf w_{K,N}(\tau)\bigr\|_{\mathcal X\times\mathcal X}^2
         & =
        \sum_{n>N}\int_0^\infty
        \eta\,\bigl\|\widehat{\mathbf w}_n(\tau,\eta)\bigr\|_{\mathbb R^2}^2\, d\eta
        +
        \sum_{n=1}^{N}\int_K^\infty
        \eta\,\bigl\|\widehat{\mathbf w}_n(\tau,\eta)\bigr\|_{\mathbb R^2}^2\, d\eta.
        \label{eq:appE-proof-tail-identity}
    \end{align}
    Thus the reconstruction error corresponds to the spectral tail in the
    \((n,\eta)\)-variables.

    \smallskip
    \noindent
    \textit{Step 2: pointwise estimate of each modal coefficient.}
    By Proposition~\ref{prop:appE-modal-stability}, there exist constants
    \(M\ge 1\) and \(\omega>0\), independent of \(n\) and \(\eta\), such that
    \begin{equation}
        \bigl\|e^{t\mathbf A_n(\eta)}\bigr\|
        \le
        M e^{-\omega t},
        \qquad t\ge 0.
        \label{eq:appE-proof-semigroup}
    \end{equation}
    Applying \eqref{eq:appE-proof-semigroup} to
    \eqref{eq:appE-proof-duhamel} gives
    \begin{align}
        \bigl\|\widehat{\mathbf w}_n(\tau,\eta)\bigr\|_{\mathbb R^2}
         & \le
        M\bigl\|\widehat{\mathbf w}_{n,0}(\eta)\bigr\|_{\mathbb R^2}
        +
        M\int_0^\tau
        e^{-\omega(\tau-\sigma)}\,|\widehat s_n(\sigma,\eta)|\, d\sigma.
        \label{eq:appE-proof-modal-bound-L1}
    \end{align}
    Using \((a+b)^2\le 2a^2+2b^2\) and then Cauchy--Schwarz in time, we obtain
    \begin{align}
        \bigl\|\widehat{\mathbf w}_n(\tau,\eta)\bigr\|_{\mathbb R^2}^2
         & \le
        2M^2\bigl\|\widehat{\mathbf w}_{n,0}(\eta)\bigr\|_{\mathbb R^2}^2
        +
        2M^2
        \left(
        \int_0^\tau
        e^{-\omega(\tau-\sigma)}|\widehat s_n(\sigma,\eta)|\, d\sigma
        \right)^2
        \notag \\
         & \le
        2M^2\bigl\|\widehat{\mathbf w}_{n,0}(\eta)\bigr\|_{\mathbb R^2}^2
        +
        2M^2
        \left(
        \int_0^\tau e^{-2\omega(\tau-\sigma)}\, d\sigma
        \right)
        \left(
        \int_0^\tau |\widehat s_n(\sigma,\eta)|^2\, d\sigma
        \right)
        \notag \\
         & \le
        2M^2\bigl\|\widehat{\mathbf w}_{n,0}(\eta)\bigr\|_{\mathbb R^2}^2
        +
        \frac{M^2}{\omega}
        \int_0^\tau |\widehat s_n(\sigma,\eta)|^2\, d\sigma
        \notag \\
         & \le
        2M^2\bigl\|\widehat{\mathbf w}_{n,0}(\eta)\bigr\|_{\mathbb R^2}^2
        +
        \frac{M^2}{\omega}
        \int_0^{\tau_f} |\widehat s_n(\sigma,\eta)|^2\, d\sigma.
        \label{eq:appE-proof-modal-bound-L2}
    \end{align}
    Therefore,
    \begin{equation}
        \sup_{0\le \tau\le \tau_f}
        \bigl\|\widehat{\mathbf w}_n(\tau,\eta)\bigr\|_{\mathbb R^2}^2
        \le
        2M^2\bigl\|\widehat{\mathbf w}_{n,0}(\eta)\bigr\|_{\mathbb R^2}^2
        +
        \frac{M^2}{\omega}
        \int_0^{\tau_f} |\widehat s_n(\sigma,\eta)|^2\, d\sigma.
        \label{eq:appE-proof-uniform-modal-bound}
    \end{equation}

    \smallskip
    \noindent
    \textit{Step 3: summation over the spectral tail.}
    Insert \eqref{eq:appE-proof-uniform-modal-bound} into
    \eqref{eq:appE-proof-tail-identity}. Since all terms are non-negative, Tonelli's theorem yields
    \begin{align}
         & \sup_{0\le \tau\le \tau_f}
        \bigl\|\mathbf w(\tau)-\mathbf w_{K,N}(\tau)\bigr\|_{\mathcal X\times\mathcal X}^2
        \notag                        \\
         & \le
        \sum_{n>N}\int_0^\infty
        \eta\,
        \sup_{0\le \tau\le \tau_f}
        \bigl\|\widehat{\mathbf w}_n(\tau,\eta)\bigr\|_{\mathbb R^2}^2
        \, d\eta
        +
        \sum_{n=1}^{N}\int_K^\infty
        \eta\,
        \sup_{0\le \tau\le \tau_f}
        \bigl\|\widehat{\mathbf w}_n(\tau,\eta)\bigr\|_{\mathbb R^2}^2
        \, d\eta
        \notag                        \\
         & \le
        2M^2
        \left[
        \sum_{n>N}\int_0^\infty
        \eta\,\bigl\|\widehat{\mathbf w}_{n,0}(\eta)\bigr\|_{\mathbb R^2}^2\, d\eta
        +
        \sum_{n=1}^{N}\int_K^\infty
        \eta\,\bigl\|\widehat{\mathbf w}_{n,0}(\eta)\bigr\|_{\mathbb R^2}^2\, d\eta
        \right]
        \notag                        \\
         & \quad
        +
        \frac{M^2}{\omega}
        \int_0^{\tau_f}
        \left[
            \sum_{n>N}\int_0^\infty
            \eta\, |\widehat s_n(\sigma,\eta)|^2\, d\eta
            +
            \sum_{n=1}^{N}\int_K^\infty
            \eta\, |\widehat s_n(\sigma,\eta)|^2\, d\eta
            \right] d\sigma.
        \label{eq:appE-proof-tail-estimate}
    \end{align}
    Hence \eqref{eq:appE-spectral-tail-bound} follows with
    \[
        C(\tau_f)\coloneqq \max\!\left\{2M^2,\frac{M^2}{\omega}\right\}.
    \]
    In particular, any larger constant depending only on \(\tau_f\) is also
    admissible.

    \smallskip
    \noindent
    \textit{Step 4: passage to the limit \(K,N\to\infty\).}
    By Parseval's identity and orthonormality of the depth basis,
    \[
        \sum_{n=1}^\infty\int_0^\infty
        \eta\,\bigl\|\widehat{\mathbf w}_{n,0}(\eta)\bigr\|_{\mathbb R^2}^2\, d\eta
        =
        \bigl\|\mathbf w(0)\bigr\|_{\mathcal X\times\mathcal X}^2
        <
        \infty,
    \]
    and similarly
    \[
        \int_0^{\tau_f}
        \sum_{n=1}^\infty\int_0^\infty
        \eta\, |\widehat s_n(\tau,\eta)|^2\, d\eta\, d\tau
        =
        \|s\|_{L^2(0,\tau_f;\mathcal X)}^2
        <
        \infty.
    \]
    Therefore the tails \(\mathcal T_0(K,N)\) and \(\mathcal T_s(K,N)\) tend to
    zero as \(K,N\to\infty\), by monotone convergence (equivalently, by the
    absolute continuity of the integral on the product measure space
    \(\mathbb N\times [0,\infty)\), with counting measure in \(n\) and
    \(\eta\,d\eta\) in the Hankel variable). Returning to
    \eqref{eq:appE-spectral-tail-bound}, we conclude that
    \[
        \sup_{0\le \tau\le \tau_f}
        \bigl\|\mathbf w(\tau)-\mathbf w_{K,N}(\tau)\bigr\|_{\mathcal X\times\mathcal X}
        \longrightarrow 0
        \qquad
        \text{as } K\to\infty,\; N\to\infty.
    \]
    This proves \eqref{eq:appE-spectral-convergence}.
\end{proof}

\begin{corollary}
    \label{cor:appE-algebraic-rate}
    Assume that there exists \(m>0\) such that
    \begin{equation}
        \sum_{n=1}^\infty \int_0^\infty
        (1+\eta^2+\mu_n^2)^m
        \eta\, \|\widehat{\mathbf{w}}_n(0,\eta)\|_{\mathbb{R}^2}^2\, d\eta
        <\infty,
        \label{eq:appE-spectral-regularity-initial}
    \end{equation}
    and similarly
    \begin{equation}
        \int_0^{\tau_f}
        \sum_{n=1}^\infty \int_0^\infty
        (1+\eta^2+\mu_n^2)^m
        \eta\, |\widehat{s}_n(\tau,\eta)|^2\, d\eta\, d\tau
        <\infty.
        \label{eq:appE-spectral-regularity-source}
    \end{equation}
    Then the truncation error satisfies the algebraic estimate
    \begin{equation}
        \sup_{0\le \tau\le \tau_f}
        \left\|
        \mathbf{w}(\tau)-\mathbf{w}_{K,N}(\tau)
        \right\|_{X\times X}
        \le
        C_m(\tau_f)\bigl(K^{-m}+N^{-m}\bigr).
        \label{eq:appE-spectral-rate}
    \end{equation}
\end{corollary}

\begin{proof}
    On the tail \(\eta\ge K\), one has
    \((1+\eta^2+\mu_n^2)^{-m}\le K^{-2m}\), while on the tail \(n>N\), one has
    \((1+\eta^2+\mu_n^2)^{-m}\le \mu_N^{-2m}\sim N^{-2m}\).
    The estimate then follows directly from
    Theorem~\ref{thm:appE-spectral-convergence}.
\end{proof}

\section{Closed-form modal source coefficients for constant optical coefficients}
\label{sec:appE-special-source-cases}

We now specialise the source to the constant-optics linear setting, namely to
the case in which the reflectivity and effective absorption coefficient are
held fixed,
\[
    R\equiv R_\star,
    \qquad
    \alpha' \equiv \alpha_\star.
\]
This is precisely the regime in which the linear pTTM reduces to a
constant-coefficient parabolic system, and it corresponds to the fixed-optics
simplification of the ULP source discussed in
Chapter~\ref{ch:theoretical-background}.
Specialising the source term inherited from the axisymmetric ULP construction (e.g., \eqref{eq:ch1-ulp_heat_source}),
one may write in solver units
\begin{equation}
    s(\tau,\xi,\zeta)
    =
    A_\star\, E(\tau)\,
    \exp\!\left\{-\frac{\xi^2}{R_0^2}\right\}\,
    \exp\{-\beta \zeta\},
    \label{eq:appE-separable-source}
\end{equation}
where \(A_\star>0\) collects all constant prefactors,
\(E(\tau)\) is the temporal envelope, and
\(\beta>0\) is the dimensionless attenuation coefficient in the depth
direction.

\subsection{Hankel transform of the Gaussian radial profile}
\label{sec:appE-gaussian-radial-transform}

The order-zero Hankel transform of the Gaussian factor is classical:
\begin{equation}
    \int_0^\infty
    \xi\,
    \exp\!\left\{-\frac{\xi^2}{R_0^2}\right\}
    J_0(\eta\xi)\, d\xi
    =
    \frac{R_0^2}{2}\,
    \exp\!\left\{-\frac{R_0^2\eta^2}{4}\right\}.
    \label{eq:appE-gaussian-hankel}
\end{equation}
Therefore the radial transform of the source remains Gaussian in transform
space.

\subsection{Depth coefficients under homogeneous Dirichlet conditions}
\label{sec:appE-depth-source-coeffs}

The Dirichlet depth coefficient is
\begin{align}
    \beta_n
     & \coloneqq
    \int_0^1 e^{-\beta\zeta}\phi_n(\zeta)\, d\zeta
    =
    \sqrt{2}\int_0^1 e^{-\beta\zeta}\sin(n\pi\zeta)\, d\zeta
    \notag       \\
     & =
    \sqrt{2}\,
    \frac{n\pi\bigl(1-(-1)^n e^{-\beta}\bigr)}
    {\beta^2+n^2\pi^2}.
    \label{eq:appE-depth-exp-decay}
\end{align}
Hence the fully transformed modal source is
\begin{equation}
    \widehat{s}_n(\tau,\eta)
    =
    \Gamma_n(\eta)\, E(\tau),
    \label{eq:appE-transformed-source}
\end{equation}
where
\begin{equation}
    \Gamma_n(\eta)
    \coloneqq
    A_\star\,
    \frac{R_0^2}{2}\,
    \exp\!\left\{-\frac{R_0^2\eta^2}{4}\right\}\,
    \beta_n.
    \label{eq:appE-gamma-n}
\end{equation}

\subsection{Closed-form modal responses}
\label{sec:appE-closed-form-modal-responses}

With zero initial data,
\(
\widehat{\mathbf{w}}_n(0,\eta)=0
\),
equation \eqref{eq:appE-duhamel} becomes
\begin{equation}
    \widehat{\mathbf{w}}_n(\tau,\eta)
    =
    \Gamma_n(\eta)
    \int_0^\tau
    e^{(\tau-\sigma)\mathbf{A}_n(\eta)}
    E(\sigma)\,\mathbf{e}_1\, d\sigma.
    \label{eq:appE-modal-source-only}
\end{equation}

\paragraph{Constant envelope.}
If \(E(\tau)\equiv 1\), then
\begin{equation}
    \widehat{\mathbf{w}}_n(\tau,\eta)
    =
    \Gamma_n(\eta)\,
    \mathbf{A}_n(\eta)^{-1}
    \bigl(
    e^{\tau \mathbf{A}_n(\eta)}-\mathbf{I}
    \bigr)\mathbf{e}_1.
    \label{eq:appE-constant-envelope}
\end{equation}

\paragraph{Exponential envelope.}
If \(E(\tau)=e^{-\sigma\tau}\) with \(\sigma>0\), then
\begin{equation}
    \widehat{\mathbf{w}}_n(\tau,\eta)
    =
    \Gamma_n(\eta)\,
    \bigl(\mathbf{A}_n(\eta)+\sigma\mathbf{I}\bigr)^{-1}
    \bigl(
    e^{\tau \mathbf{A}_n(\eta)}-e^{-\sigma\tau}\mathbf{I}
    \bigr)\mathbf{e}_1.
    \label{eq:appE-exponential-envelope}
\end{equation}

\paragraph{Gaussian envelope.}
If
\begin{equation}
    E(\tau)=\exp\!\bigl\{-\kappa(\tau-\tau_c)^2\bigr\},
    \qquad
    \kappa>0,
    \label{eq:appE-gaussian-envelope}
\end{equation}
and
\(
\mathbf{A}_n(\eta)=\mathbf{P}_n(\eta)\mathbf{\Lambda}_n(\eta)\mathbf{P}_n(\eta)^{-1}
\)
with
\[
    \mathbf{\Lambda}_n(\eta)
    =
    \operatorname{diag}\bigl(\lambda_{+,n}(\eta),\lambda_{-,n}(\eta)\bigr),
\]
then
\begin{equation}
    \widehat{\mathbf{w}}_n(\tau,\eta)
    =
    \Gamma_n(\eta)\,
    \mathbf{P}_n(\eta)
    \begin{bmatrix}
        \mathcal{I}_{\lambda_{+,n}(\eta)}(\tau) & 0                                       \\
        0                                       & \mathcal{I}_{\lambda_{-,n}(\eta)}(\tau)
    \end{bmatrix}
    \mathbf{P}_n(\eta)^{-1}\mathbf{e}_1,
    \label{eq:appE-gaussian-modal-solution}
\end{equation}
where
\begin{align}
    \mathcal{I}_{\lambda}(\tau)
     & \coloneqq
    \int_0^\tau
    e^{\lambda(\tau-\sigma)}
    e^{-\kappa(\sigma-\tau_c)^2}\, d\sigma
    \notag       \\
     & =
    \frac{\sqrt{\pi}}{2\sqrt{\kappa}}\,
    \exp\!\left\{
    \lambda(\tau-\tau_c)+\frac{\lambda^2}{4\kappa}
    \right\}
    \left[
        \operatorname{erf}\!\left(
        \sqrt{\kappa}(\tau-\tau_c)-\frac{\lambda}{2\sqrt{\kappa}}
        \right)
        -
        \operatorname{erf}\!\left(
        -\sqrt{\kappa}\,\tau_c-\frac{\lambda}{2\sqrt{\kappa}}
        \right)
        \right].
    \label{eq:appE-gaussian-convolution}
\end{align}

\paragraph{Lorentzian envelope.}
In the notation of Chapter~\ref{ch:theoretical-background}, the Lorentzian pulse centered at \(\tau_c\)
with half-width at half-maximum \(\gamma>0\) is
\begin{equation}
    E(\tau)
    =
    \frac{1}{\pi\gamma}
    \left[
        1+\left(\frac{\tau-\tau_c}{\gamma}\right)^2
        \right]^{-1},
    \label{eq:appE-lorentzian-envelope}
\end{equation}
which reduces to \eqref{eq:ch1-lorentzian} by letting \(\tau_p \coloneqq 2\gamma\) denote the full width at half maximum,
and after passage to scaled time.

Assuming again that
\(
\mathbf{A}_n(\eta)=\mathbf{P}_n(\eta)\mathbf{\Lambda}_n(\eta)\mathbf{P}_n(\eta)^{-1}
\)
with
\[
    \mathbf{\Lambda}_n(\eta)
    =
    \operatorname{diag}\bigl(\lambda_{+,n}(\eta),\lambda_{-,n}(\eta)\bigr),
\]
the modal response takes the form
\begin{equation}
    \widehat{\mathbf{w}}_n(\tau,\eta)
    =
    \Gamma_n(\eta)\,
    \mathbf{P}_n(\eta)
    \begin{bmatrix}
        \mathcal{L}_{\lambda_{+,n}(\eta)}(\tau) & 0                                       \\
        0                                       & \mathcal{L}_{\lambda_{-,n}(\eta)}(\tau)
    \end{bmatrix}
    \mathbf{P}_n(\eta)^{-1}\mathbf{e}_1,
    \label{eq:appE-lorentzian-modal-solution}
\end{equation}
where
\begin{align}
    \mathcal{L}_{\lambda}(\tau)
     & \coloneqq
    \int_0^\tau
    e^{\lambda(\tau-\sigma)}
    \frac{1}{\pi\gamma}
    \left[
        1+\left(\frac{\sigma-\tau_c}{\gamma}\right)^2
        \right]^{-1}
    d\sigma
    \notag            \\
     & =
    \frac{e^{\lambda(\tau-\tau_c)}}{2\pi i}
    \Bigl[
        e^{-i\lambda\gamma}
        \Bigl(
        \operatorname{Ei}\!\bigl(-\lambda(\tau-\tau_c)+i\lambda\gamma\bigr)
        -
        \operatorname{Ei}\!\bigl(\lambda\tau_c+i\lambda\gamma\bigr)
        \Bigr)
    \notag            \\
     & \hspace{4.5em}
        -
        e^{i\lambda\gamma}
        \Bigl(
        \operatorname{Ei}\!\bigl(-\lambda(\tau-\tau_c)-i\lambda\gamma\bigr)
        -
        \operatorname{Ei}\!\bigl(\lambda\tau_c-i\lambda\gamma\bigr)
        \Bigr)
        \Bigr].
    \label{eq:appE-lorentzian-convolution}
\end{align}
Here \(\operatorname{Ei}\) denotes the exponential-integral function on its
principal branch. The formula is obtained by the substitution
\(
u=(\sigma-\tau_c)/\gamma
\)
and the partial-fraction identity
\[
    \frac{1}{1+u^2}
    =
    \frac{1}{2i}
    \left(
    \frac{1}{u-i}-\frac{1}{u+i}
    \right).
\]

As a consistency check, in the limit \(\lambda\to 0\),
\eqref{eq:appE-lorentzian-convolution} reduces to the elementary primitive
\begin{equation}
    \mathcal{L}_{0}(\tau)
    =
    \frac{1}{\pi}
    \left[
        \arctan\!\left(\frac{\tau-\tau_c}{\gamma}\right)
        +
        \arctan\!\left(\frac{\tau_c}{\gamma}\right)
        \right].
    \label{eq:appE-lorentzian-zero-mode-check}
\end{equation}

\paragraph{Remark on double-pulse forcing.}
Because the modal problem \eqref{eq:appE-modal-source-only} is linear, a
double-pulse envelope of the form
\[
    E_{\mathrm{DP}}(\tau)
    =
    E_1(\tau-\tau_{c,1})+E_2(\tau-\tau_{c,2})
\]
is treated by superposition: the corresponding modal response is the sum of
the two single-pulse responses computed from
\eqref{eq:appE-gaussian-modal-solution} or
\eqref{eq:appE-lorentzian-modal-solution}, respectively.

These formulae provide closed-form modal responses for the most common
constant-optics heat-source specialisations relevant to the linear solver.

\section{Remarks on alternative depth conditions and finite-radius truncation}
\label{sec:appE-extensions}

\begin{remark}[Homogeneous Neumann depth conditions]
    \label{rem:appE-neumann-case}
    If homogeneous Neumann conditions are imposed in the depth variable, then the
    appropriate basis is the cosine family
    \[
        \phi_n^{(N)}(\zeta)=
        \begin{cases}
            1,                       & n=0,    \\[2pt]
            \sqrt{2}\cos(n\pi\zeta), & n\ge 1,
        \end{cases}
        \qquad
        \mu_n=n\pi.
    \]
    The modal derivation remains unchanged after replacing the sine basis
    \eqref{eq:appE-depth-basis} by this cosine basis. In that case the depth
    transform is the one most closely associated with a cosine-transform or DCT
    implementation \cite{Rogers_2012}.
\end{remark}

\begin{remark}[Finite radial domain]
    If the radial variable is truncated to \(\xi\in[0,\xi_{\max}]\) rather than
    \([0,\infty)\), then one should replace the infinite Hankel transform by a
    finite Hankel transform or a generalized finite Hankel transform. This is the
    natural setting for the constructions developed by Cinelli
    \cite{Cinelli_1965}. Since the present appendix is concerned with the
    continuous infinite-radius axisymmetric formulation, we do not pursue that
    variant here.
\end{remark}

\begin{remark}[Relation to the proposed numerical pipeline]
    The appendix establishes the continuous spectral structure underlying the
    proposed transform-based axisymmetric solver. In particular, it justifies the
    radial use of an order-zero Hankel transform and the depth use of a sine or
    cosine spectral basis according to the depth boundary conditions.
    A pseudo-spectral implementation based on operator splitting and a discrete or
    quasi-discrete Hankel transform may therefore be interpreted as a consistent
    discretization of the continuous modal system
    \eqref{eq:appE-modal-ode}. This observation is consistent with transform-based
    axisymmetric heat-conduction solvers already used in the literature
    \cite{Rogers_2012,Kunga_2021}.
\end{remark}
